%=============================================================================!
% 7.7  STEPS THAT KEEP AREA                                                   !
%=============================================================================!
\section{Steps that keep area}
\label{sec:ha-symplectic}

A computer follows a motion in steps. A method of computing motion takes
the state $(q, p)$ at one instant and returns the state $(q', p')$ a time
$\dt$ later, so it is a map of the phase plane into itself, a step map,
and like the flow it has a Jacobian determinant. The exact flow over $\dt$
keeps both the energy and areas. This section shows, for the spring and
the two simplest methods, that a step map which inflates areas makes the
energy grow without limit, while the area-preserving method holds the energy within a band for
the stable step used here.

\subsection*{The forward Euler method}

The simplest step takes the rates at the start of the step as if they held throughout it, so that the change over the step is the
rate times $\dt$, the secant of \cref{sec:mo-velocity} read backwards:
\begin{equation*}
   q' = q + \dt\,\frac{\partial H}{\partial p}(q, p) , \qquad
   p' = p - \dt\,\frac{\partial H}{\partial q}(q, p) .
\end{equation*}
This is the forward Euler method. For the block of \cref{sec:os-circle},
\SI{0.5}{\kilogram} on a spring of \SI{50}{\newton\per\metre}, it reads
$q' = q + \dt\,p/m$ and $p' = p - \dt\,kq$. Its energy after the step is
\begin{align*}
   H' &= \frac{p'^2}{2m} + \frac12 kq'^2
   = \frac{(p - \dt\,kq)^2}{2m} + \frac12 k\Bigl(q + \frac{\dt\,p}{m}
   \Bigr)^2\\
   &= \frac{p^2}{2m} - \frac{\dt\,kqp}{m} + \frac{\dt^2k^2q^2}{2m}
   + \frac12 kq^2 + \frac{\dt\,kqp}{m} + \frac{\dt^2kp^2}{2m^2}
   && \text{expanding both squares}\\
   &= \frac{p^2}{2m} + \frac12 kq^2
   + \frac{k}{m}\,\dt^2\Bigl(\frac12 kq^2 + \frac{p^2}{2m}\Bigr)
   && \text{collecting terms}\\
   &= \bigl(1 + \omega^2\dt^2\bigr)\,H
   && \text{with $\omega^2 = k/m$.}
\end{align*}
In the third line the two terms in $qp$ have cancelled and $k\dt^2/m$ has
been taken out of the two terms in $\dt^2$. Every step multiplies the
energy by $1 + \omega^2\dt^2$, whatever the state, so the computed motion spirals outwards in the phase plane
(\cref{fig:ha-maps}). With $\omega = \SI{10}{\radian\per\second}$ and $\dt
= \SI{0.02}{\second}$, so that $\omega\dt = 0.2$, the factor is 1.04, and
after 50 steps, one second, the energy is 7.107 times its starting value;
after \SI{10}{\second} it is $3.3\times10^8$ times. A shorter step only
slows the growth: the energy doubles after the $n$ steps for which $(1 +
\omega^2\dt^2)^n = 2$. Writing $1 + \omega^2\dt^2 = \mathrm{e}^{\ln(1 +
\omega^2\dt^2)}$ and $2 = \mathrm{e}^{\ln 2}$ (\cref{sec:ne-drag}), this
is $n\ln(1 + \omega^2\dt^2) = \ln 2$, so $n = \ln 2/\ln(1 +
\omega^2\dt^2)$, which is 17.7 at $\omega\dt = 0.2$, 69.7 at 0.1 and 6932
at 0.01, about 11 periods of $2\pi/(\omega\dt) = 628$ steps each.

The step map also stretches areas. Its Jacobian matrix, from the
derivatives of $q'$ and $p'$ with respect to $q$ and $p$, and its
determinant are
\begin{equation*}
   \vb{J} = \begin{pmatrix} 1 & \dt/m\\ -k\,\dt & 1 \end{pmatrix} ,
   \qquad
   \det\vb{J} = 1\times1 - \frac{\dt}{m}\times(-k\,\dt)
   = 1 + \omega^2\dt^2 ,
\end{equation*}
the same factor: every step inflates every area by it.

\subsection*{The symplectic Euler method}

A small change to the order of the updates gives a map that keeps area.
The momentum is moved first, with the force at the old position, and then
the position, with the new momentum:
\begin{equation*}
   p' = p - \dt\,\frac{\partial H}{\partial q}(q) , \qquad
   q' = q + \dt\,\frac{\partial H}{\partial p}(p') ,
\end{equation*}
for a Hamiltonian $H = K(p) + U(q)$ whose kinetic energy depends only on
the momentum and potential energy only on the position, as for atoms in
Cartesian coordinates. For the spring, $p' = p - \dt\,kq$, and putting
this into $q' = q + \dt\,p'/m$ gives $q' = (1 - \omega^2\dt^2)\,q + \dt\,
p/m$. So
\begin{equation*}
   \vb{J} = \begin{pmatrix} 1 - \omega^2\dt^2 & \dt/m\\ -k\,\dt & 1
   \end{pmatrix} ,
   \qquad
   \det\vb{J} = 1 - \omega^2\dt^2 - \frac{\dt}{m}\times(-k\,\dt) = 1 .
\end{equation*}
The same holds for any $H = K(p) + U(q)$. The step is two maps in turn: a
kick, $(q, p) \to (q, p - \dt\,U'(q))$, which changes only the momentum,
by an amount that depends only on the position, and then a drift, $(q, p')
\to (q + \dt\,K'(p'), p')$, which changes only the position, by an amount
that depends only on the momentum. Their Jacobian matrices are
\begin{equation*}
   \begin{pmatrix} 1 & 0\\ -\dt\,U''(q) & 1 \end{pmatrix}
   \qquad\text{and}\qquad
   \begin{pmatrix} 1 & \dt\,K''(p')\\ 0 & 1 \end{pmatrix} ,
\end{equation*}
each with determinant $1\times1 - 0 = 1$: each is a shear. A small square
taken through the kick keeps its area, and the parallelogram it becomes
keeps its area through the drift, so the whole step keeps area. This is
the symplectic Euler method; a step map that keeps area in the phase
plane is called symplectic. With more coordinates, keeping volume is
necessary but the full condition is stronger, and Chapter~9 states
it\cite{LeimkuhlerMatthews2015}. The function \texttt{symplectic\_euler\_step}
of \texttt{mdlab} takes one step:

\begin{code}
def symplectic_euler_step(
    dh_dq: Callable[[NDArray, NDArray], NDArray],
    dh_dp: Callable[[NDArray, NDArray], NDArray],
    q: ArrayLike,
    p: ArrayLike,
    dt: float,
) -> tuple[NDArray, NDArray]:
    q = np.asarray(q, dtype=float)
    p = np.asarray(p, dtype=float)
    p_new = p - dt * dh_dq(q, p)
    return q + dt * dh_dp(q, p_new), p_new
\end{code}

\noindent The two derivatives of $H$ are passed in as functions, called
with the state as it stands: the kick with the old $q$, the drift with the
new momentum \texttt{p\_new}. For the pendulum of \cref{sec:ha-equations},
whose step is not linear, the determinant computed by central differences
with $\dt = \SI{0.05}{\second}$ differs from 1 by less than $10^{-9}$ at
each of the states tried.

\begin{figure}[htb]
\centering
\includegraphics{ch07_hamilton/maps}
\caption{The block of \SI{0.5}{\kilogram} on a spring of
\SI{50}{\newton\per\metre} ($\omega = \SI{10}{\radian\per\second}$),
released from rest at \SI{4}{\centi\metre}, stepped with $\dt =
\SI{0.02}{\second}$ by the forward Euler method (oxblood) and the
symplectic Euler method (green). (a) The states after each step for
\SI{1}{\second}, with the true ellipse (dashed), which the green points
nearly hide. (b) The energy over \SI{10}{\second} as a
multiple of its starting value, on a logarithmic axis; the forward Euler
energy leaves the axis, at 1000 times its starting value, after
\SI{3.5}{\second}.}
\label{fig:ha-maps}
\end{figure}

With the same step of \SI{0.02}{\second}, the symplectic Euler method keeps
the block's energy between 0.9091 and 1.1111 times its starting value,
over \SI{10}{\second} and over \SI{2000}{\second} alike
(\cref{fig:ha-maps}b): its error oscillates and does not grow. Halving the
step to \SI{0.01}{\second} narrows the band to between 0.9524 and 1.0526.
Chapter~9 explains why: a step map that keeps area keeps a slightly
different energy, exactly for the spring and very nearly in
general\cite{LeimkuhlerMatthews2015}, and so stays close to the true one.
It also finds methods whose band is much narrower for the same step.

\begin{practice}
The total energy of an isolated system is the first diagnostic of a
method of computing motion and of its step. A steady growth or fall shows
a method that does not keep area, as here, a step too long for the
method, or forces that are not the slopes of a single potential energy
(\cref{sec:en-drift}); for smooth forces and sufficiently short stable steps, the symplectic
method shown here has bounded fluctuations that shrink as the step is
shortened. Chapter~9 sizes the step from these
fluctuations.
\end{practice}

\begin{keyidea}
A step map keeps area when its Jacobian determinant is 1. The forward Euler
method multiplies the spring's energy and every area by $1 + \omega^2
\dt^2$ at each step; the symplectic Euler method, a kick then a drift, keeps
area and holds the energy within a band.
\end{keyidea}

\notebook{07}{step the spring with both methods and watch a square of states}

\begin{selfcheck}
With $\omega\dt = 0.05$, by what factor does the forward Euler method
multiply the energy of a spring in 100 steps?
\end{selfcheck}
